%##%\FINALFIX{\newpage}
\subsection{nml\_reference}
\label{SEC:nml_reference}
\inp*[nmlReference]{}[nml\_reference][]

This \hyperref[SEC:User_Input]{namelist block}
is used to specify the reference conditions for the simulation.
%
These conditions are used within the flow solver to nondimensionalize
the internal variables, and in most cases to initialize the
volume solution.
%
They are also written to the screen during code execution.

\begin{nmlblocklst}[
,label={LST:nml_reference}
,caption={Namelist nml\_reference.}
]
&nml_reference
  Mach         = 0.4
  Reynolds     = -1.0      ! rho*U*L/mu, L=1.0[m]
  Gamma        = 1.4
  Gas_Constant = 287.05    ! [J/(kg*K)] = [m^2/(s^2*K)]
  Density      = -1.0      ! [kg/m^3]
  Temperature  = 300.0     ! [K]
  Pressure     = 100000.0  ! [Pa]
  AoA_Alpha    = 0.0       ! [deg]
  AoA_Beta     = 0.0       ! [deg]
  Prandtl      = 0.72
  Schmidt      = 0.72
  Suth_Muref   = 1.716e-5  ! [Pa*s] = [N*s/m^2]
  Suth_Tref    = 273.0     ! [K]
  Suth_Sref    = 110.4     ! [K]
/
\end{nmlblocklst}

\NOTE
Certain initialization options replace the reference conditions
by preset values to ensure repeatability of these benchmark cases.
%
See the \inp[nmlInitialization::Initialization]{Initialization}
input for details.

%==========================================================================

\begin{nmldescitem}
Mach = 0.4
\end{nmldescitem}
\inp*[nmlReference::Mach]{}[Mach][]

\begin{nmldescbody}
Reference Mach number.
Must be non-negative.
\end{nmldescbody}

%==========================================================================

\begin{nmldescitem}
Reynolds = -1.0  ! rho*U*L/mu, L=1.0[m]
\end{nmldescitem}
\inp*[nmlReference::Reynolds]{}[Reynolds][]

\begin{nmldescbody}
Reference Reynolds number, $Re = \rho U L / \mu$, based on
reference conditions.
%
Here,
$\rho$ is the reference density;
$U$ is the velocity magnitude computed from
  the reference Mach number and
  speed of sound at the reference temperature;
$L$ is the reference length of 1.0 meter; and 
$\mu$ is the viscosity computed from
  Sutherland's viscosity law using \eref{EQN:Viscosity_Sutherland}
  at the reference temperature.

How the Reynolds number is used depends on the input value:
\\
\begin{tabular}{ @{\qquad} r @{ \quad } p{0.85\linewidth} @{} }
%\begin{longtable}{ @{\qquad} r @{ \quad } p{0.85\linewidth} @{} }
  $\leq0$
       & The Reynolds number will be computed from the
         other reference flow properties.
         %
         This yields \scinum{8.7319}{6} for the default values shown above,
         but will change if you modify any of the other reference conditions.
         \\
  $>0$ & The Reynolds number will be used with the
         other reference conditions you specify
         to recalibrate the $\mu_0$ constant
         (\inp[nmlReference::SuthMuref]{Suth\_Muref})
         in \key[key:SutherlandsViscosityLaw]{Sutherland's viscosity law}[=],
         see \eref{EQN:Viscosity_Sutherland}.
         %
         This can be used to prescribe an artificially lower Reynolds number
         that permits better resolution of turbulent scales with
         large-eddy simulation.
         %
         It cannot be used to compute other unspecified conditions such as
         density, pressure, or temperature.
%\end{longtable}
\end{tabular}
\end{nmldescbody}

%==========================================================================

\FINALFIX{\newpage}
\begin{nmldescitem}
Gamma = 1.4
\end{nmldescitem}
\inp*[nmlReference::Gamma]{}[Gamma][]

\begin{nmldescbody}
Reference $\gamma$ (ratio of specific heats).
%
If the value specified is negative, it will be reset to the default
of 1.4 (the typical value for air).
\end{nmldescbody}

%==========================================================================

\begin{nmldescitem}
Gas_Constant = 287.05  ! [J/(kg*K)] = [m^2/(s^2*K)]
\end{nmldescitem}
\inp*[nmlReference::GasConstant]{}[Gas\_Constant][]

\begin{nmldescbody}
Reference specific gas constant (units of \unit{\J\per\kg\per\K}).
%
If the value specified is negative, it will be reset to the default
of 287.05 (the typical value for air).
 
\noindent\NOTE
In the code, the default value actually has more precision since it is computed
as $\texttt{Gas\_Constant} = R_\textrm{univ} / M_\textrm{air}$,
where $R_\textrm{univ}$ is the universal gas constant that is hard-coded as a
Fortran constant with the value
$8314.5\ \unit{\J\per\K\per\kmol}$,
and $M_\textrm{air}$ is the
molar mass of air which is also a hard-coded Fortran constant with the value
$28.9651159\ \unit{\kg\per\kmol}$.
\end{nmldescbody}

%==========================================================================

\begin{nmldescitem}
Density     = -1.0      ! [kg/m^3]
Temperature = 300.0     ! [K]
Pressure    = 100000.0  ! [Pa]
\end{nmldescitem}
\inp*[nmlReference::Density]{}[Density][]
\inp*[nmlReference::Temperature]{}[Temperature][]
\inp*[nmlReference::Pressure]{}[Pressure][]

\begin{nmldescbody}
Reference static
density $\rho$ (\unitDens{}),
temperature $T$ (\unitTemp{}), and 
pressure $p$ (\unitPres{}).

\NOTE
If all three inputs are positive, a consistency check will be performed.
%
If one of these three inputs is negative, it will be computed
from the other two using the equation of state, \eref{EQN:EqOfState}.
%
So, for the default values shown above,
density will evaluate to 1.160833 \unitDens{}.
\end{nmldescbody}

%==========================================================================

\begin{nmldescitem}
AoA_Alpha = 0.0  ! [deg]
AoA_Beta  = 0.0  ! [deg]
\end{nmldescitem}
\inp*[nmlReference::AoAAlpha]{}[AoA\_Alpha][]
\inp*[nmlReference::AoABeta]{}[AoA\_Beta][]

\begin{nmldescbody}
Reference flow angles (in degrees):
$\alpha$ is the
\key[key:AngleOfAttack]{angle of attack}[flow angles!angle of attack]
measured in the $x$-$y$ plane, and
$\beta$  is the
\key[key:AngleOfSideslip]{angle of sideslip}[flow angles!angle of sideslip]
measured in the $x$-$z$ plane.
%
These flow angles are defined by
\eref{EQN:FlowAngleVelX,EQN:FlowAngleVelY,EQN:FlowAngleVelZ} and
\fref{FIG:FlowAngle} on \pref{\pageref{FIG:FlowAngle}}.
\end{nmldescbody}

%==========================================================================

\begin{nmldescitem}
Prandtl = 0.72
\end{nmldescitem}
\inp*[nmlReference::Prandtl]{}[Prandtl][]

\begin{nmldescbody}
Prandtl number of the flow.
\end{nmldescbody}

%==========================================================================

\begin{nmldescitem}
Schmidt = 0.72
\end{nmldescitem}
\inp*[nmlReference::Schmidt]{}[Schmidt][]

\begin{nmldescbody}
Schmidt number of the flow.
\end{nmldescbody}

%==========================================================================

\begin{nmldescitem}
Suth_Muref = 1.716e-5  ! [Pa*s] = [N*s/m^2]
Suth_Tref  = 273.0     ! [K]
Suth_Sref  = 110.4     ! [K]
\end{nmldescitem}
\inp*[nmlReference::SuthMuref]{}[Suth\_Muref][]
\inp*[nmlReference::SuthTref]{}[Suth\_Tref][]
\inp*[nmlReference::SuthSref]{}[Suth\_Sref][]

\begin{nmldescbody}
These inputs correspond to the constants
  $\mu_0$,
  $T_0$, and
  $S_0$
in \key[key:SutherlandsViscosityLaw]{Sutherland's viscosity law}[=],
\eref{EQN:Viscosity_Sutherland}.

\noindent
\NOTE
If the reference Reynolds number specified is positive,
it will be used to recalibrate the $\mu_0$ constant.
See \inp[nmlReference::Reynolds]{Reynolds} for details.
\end{nmldescbody}

%==========================================================================

